\section{Spin structures and the signs of covers}\label{sec:spin}

Let $C$ be a closed Riemann surface of genus $g\ge2$, and write $K_C$
for its canonical bundle. A theta characteristic is a holomorphic line
bundle $L$ with a specified isomorphism $L^{\otimes2}\simeq K_C$.
Given two theta characteristics $L_1,L_2$, put $E=L_1\oplus L_2$.
The two square maps define a nondegenerate holomorphic symmetric
pairing $\gamma:E\otimes E\to K_C$, with the summands orthogonal.
Fix a Hermitian metric $h$ on $E$. Applying its conjugate-linear
identification $E^*\to E$ to $\gamma(s,\cdot)$ and conjugating the
one-form factor gives a conjugate-linear map
$R:E\to\Lambda^{0,1}T^*C\otimes E$. Our endpoint operator is
\[
                 D=\bar\partial_E+R .
\]
A real Cauchy--Riemann operator with the same symbol differs from $D$
by a real bundle homomorphism of order zero.

Give the total space of $E$ the almost-complex structure for which
the projection is holomorphic, each fiber has its complex linear
structure, and the graph equation over a holomorphic map
$f:\Sigma\to C$ is $D_f s=0$. Here $D_f$ is the pullback real
Cauchy--Riemann operator, with $df$ applied to the differential
factor of $R$. In holomorphic bundle coordinates its horizontal lift
of $\partial_{\bar z}$ has vertical $(1,0)$ part $-R_{\bar z}(v)$
at $(z,v)$. The ambient first Chern number on the zero section is
$2-2g+\deg E=0$.

\begin{lemma}\label{lem:spin-gap}
For $D=\bar\partial_E+R$, the zero section is the only compact
irreducible curve in the total space. With metrics pulled back from
$C$, there exists $b>0$ such that on every finite \'etale cover $f$
the pulled-back operator $D_f$ has index zero and satisfies
\[
                       \|D_f s\|_{L^2}\ge b\|s\|_{L^2}.
\]
The same inequality, with $b/2$, holds simultaneously on all such
covers after a zeroth-order perturbation of pointwise operator norm
less than $b/2$.
\end{lemma}

\begin{proof}
For a compact curve projecting nonconstantly to $C$, let $f:\Sigma\to C$
be the projection of its normalization and let $s$ be its vertical
coordinate. The form $\alpha=df^*\gamma(s,s)$ on $\Sigma$ is globally
defined, where $df^*$ denotes pullback on the one-form factor.
Holomorphicity of $\gamma$ and symmetry give
$\bar\partial\alpha=2\,df^*\gamma(\bar\partial s,s)$, with exterior
multiplication in the differential factors. In coordinates $u$ on
$\Sigma$, $z$ on $C$, and $\gamma=G\,dz$, the equation for the graph
therefore gives
\[
  \bar\partial\alpha
    =-2\,|G(s,\cdot)|_{h^*}^{2}\,|\partial_u f|^2\,d\bar u\wedge du .
\]
Stokes' theorem implies $s=0$ on the unramified locus and hence on
$\Sigma$. A compact curve in a single complex vector space is constant.

For an \'etale cover and an arbitrary smooth section, the mixed
$L^2$ inner product of $\bar\partial s$ with $R_f s$ is a fixed
constant multiple of $\int_\Sigma\bar\partial\alpha$. Indeed, in a
conformal coordinate the factor in the cotangent metric cancels the
factor in the volume form. Thus for every real $a$,
\begin{equation}\label{eq:spin-norm}
 \|\bar\partial s+a R_f s\|_2^2
   =\|\bar\partial s\|_2^2+a^2\|R_f s\|_2^2 .
\end{equation}
The least singular value of $R$ has a positive minimum $b$ on $C$
and retains that bound under \'etale pullback. Riemann--Hurwitz
gives $g(\Sigma)-1=(\deg f)(g-1)$ on each connected cover.
Riemann--Roch then gives complex Dolbeault index zero for $f^*E$,
and hence real index zero for $D_f$. The perturbation bound follows
from the triangle inequality.
\end{proof}

Let $P_f$ be the real permutation local system on the sheets of $f$.
The evaluation of $w_1(P_f)$ on a loop is the parity of its
monodromy permutation. Write $\ell_w=L_2L_1^{-1}$ for the flat
sign line determined by $w\in H^1(C;\ZZ/2)$, using its specified
square trivialization, and set
\[
            \epsilon=(-1)^{h^0(C,L_1)+h^0(C,L_2)} .
\]
We orient the determinant line of a real Cauchy--Riemann operator
by continuation of the complex orientation through the affine
operator space. At an invertible operator its comparison with the
canonical orientation of the determinant of the zero spaces is a
sign.

\begin{proposition}\label{prop:spin-sign}
For any finite \'etale cover $f$ of degree $m$, the determinant sign of
$D_f$ is
\begin{equation}\label{eq:spin-sign}
        \sigma(f)=\epsilon^m
        (-1)^{\langle w\smile w_1(P_f),[C]\rangle}.
\end{equation}
For each $w\ne0$ there are pairs with this $w$ and either value of
$\epsilon$. A pair of identical theta characteristics gives $w=0$
and $\epsilon=1$. These pairs can all be realized on the same
smooth complex rank-two bundle over $C$.
\end{proposition}

\begin{proof}
First follow $\bar\partial+aR_f$ from $a=0$ to $a=1$ on the source.
Identify the complex cokernel of $\bar\partial$ with
$H^0(f^*E)^*$ by Serre duality and the pulled-back symmetric pairing,
using $f^*K_C=K_\Sigma$. The map from the kernel to the cokernel
induced by the derivative $R_f$ is a conjugate-linear isomorphism.
Indeed, its Serre pairing on two holomorphic sections is a nonzero
constant times the integrated Hermitian product of their images under
the symmetric isomorphism. The corresponding real determinant has
sign $(-1)^{h^0(f^*E)}$. Splitting off complex complements gives a
finite-dimensional reduced map $a[R_f]+O(a^2)$ with that sign for
small $a>0$. Equation~\eqref{eq:spin-norm} extends it to $a=1$.
This is the two-line version of the spin determinant computation
of Lee and Parker \cite[Theorem~3.1]{lp}.

Let $q_L$ be the spin quadratic refinement of the mod-two intersection
form on $H_1(C;\ZZ/2)$. Spin parity equals its Arf invariant, and
twisting by $w$ adds intersection with $x=\operatorname{PD}(w)$
to $q_L$. The Arf formula in a symplectic basis therefore gives
\[
       h^0(L\ell_w)+h^0(L)=q_L(x)\pmod2
\]
\cite{atiyah,johnson}. For $w\ne0$, choose a simple nonseparating
circle $c$ representing $x$. Its preimage consists of disjoint
circles $c_i$ covering $c$ with degrees $n_i$. For a simple circle
the quadratic refinement is one plus the evaluation of the spin
double cover on the tangentially framed circle
\cite[Theorems~1A and~3A]{johnson}. The tangent framing pulls back
along the degree-$n_i$ cover. Hence
\[
       q_{f^*L}([c_i])=n_i q_L(x)+n_i-1\pmod2 .
\]
The union represents $\operatorname{PD}(f^*w)$ and all pairwise
intersection terms vanish. Summing and using $\sum_i n_i=m$ gives
$m q_L(x)$ plus the parity of the permutation around $c$.
This proves~\eqref{eq:spin-sign} for $L=L_1$. For $w=0$ the two
parities on the cover coincide. On a disconnected cover the
determinants multiply.

Twisting both theta characteristics by a sign line whose class
evaluates nontrivially on $x$ reverses $q_{L_1}(x)$ and preserves
$w$. This realizes both signs for $w\ne0$. All the resulting bundles
have rank two and degree $2g-2$, which determine their smooth
complex bundle isomorphism class on a surface. After any such
identification the operators are joined by a smooth path of
zeroth-order perturbations. Complex bundle isomorphisms preserve
the determinant convention.
\end{proof}

The sign formula also gives the determinant of the twist $D^U$ by
any flat real orthogonal local system $U$ of rank $r$. The pairing
on $U$ and the argument for~\eqref{eq:spin-norm} give invertibility.
For a theta characteristic $L$ the orthogonal mod-two index formula
is
\[
 h^0(L\otimes U_{\CC})\equiv
 r h^0(L)+q_L(\operatorname{PD}(w_1(U)))
              +\langle w_2(U),[C]\rangle \pmod2
\]
\cite[Lemma~5.8]{bspectral}. Adding this equality for $L=L_1,L_2$
cancels the second Stiefel--Whitney terms. Thus
\begin{equation}\label{eq:spin-twist}
        \operatorname{sign}(D^U)
            =\epsilon^r(-1)^{\langle w\smile w_1(U),[C]\rangle}.
\end{equation}

Finally, for a connected stable map of covering degree $m>0$ in an
endpoint tube, let $d_i$ be the positive degrees of its nonconstant
normalized components. Each has genus at least $1+d_i(g-1)$.
The arithmetic genus is the sum of the normalized genera and the
first Betti number of the dual graph. Its minimum $1+m(g-1)$
therefore requires one nonconstant component, \'etale over $C$,
and a tree of components of which every contracted one is rational.
Any nonempty contracted subtree has an unstable leaf for unmarked
maps. The maps at the minimum are precisely the connected \'etale
covers. The derivative $df$ identifies the tangent deformation
complex in the target curve with the source deformation complex
with its complex orientation. The remaining normal operator $D_f$
is invertible. The regular local contribution of this map in the
fixed target is consequently $\sigma(f)/|\operatorname{Aut}_C(f)|$.
